\documentclass[11pt,a4paper]{article}
\usepackage[utf8]{inputenc}
\usepackage[T1]{fontenc}
\usepackage[spanish,es-nodecimaldot]{babel}
\usepackage{amsmath,amssymb}
\usepackage[round]{natbib}
\usepackage{booktabs}
\usepackage{graphicx}
\usepackage[margin=2.5cm]{geometry}
\usepackage{hyperref}
\usepackage{xcolor}
\hypersetup{colorlinks=true,linkcolor=blue!50!black,citecolor=blue!50!black,urlcolor=blue!50!black}
\usepackage{caption}
\captionsetup{font=small,labelfont=bf}

\title{\bfseries La consulta no puede ver la pregunta\\
\large El alcance de una convolución corta decide qué parte de una consulta condiciona\\
la recuperación, y lo que queda afuera se vuelve una respuesta equivocada y confiada}

\author{Maximiliano Speranza\\
\small Investigador independiente, Buenos Aires, Argentina\\
\small \href{https://orcid.org/0009-0005-0413-8554}{ORCID 0009-0005-0413-8554}\\
\small \texttt{maximiliano.speranza@gmail.com}}

\date{Septiembre de 2026}

\begin{document}
\maketitle

\begin{abstract}
\noindent
Los modelos de atención lineal y de espacio de estados arman su consulta con una convolución causal
corta. El tamaño de núcleo $4$ es el valor por omisión de hecho, así que la consulta de una capa
resulta ser función del token actual y de tres anteriores. Mostramos que ese alcance no es un detalle
de suavizado local sino un límite duro sobre \emph{qué parte de una pregunta puede condicionar la
recuperación}, y que lo que cae afuera no se degrada de a poco. Desaparece exacto, y el modelo
contesta con confianza a partir de la parte que todavía alcanza a ver.

Medimos el efecto en un modelo de lenguaje chico con un archivo persistente co-entrenado que se
consulta entre secuencias. La sensibilidad de la recuperación a un token de la consulta, medida como
el movimiento en variación total de la distribución de lectura cuando ese token se reemplaza, vale
$\mathbf{0{,}000000}$ apenas el token pasa el alcance de la convolución, y el corte se mueve con el
núcleo. Con núcleo $3$ (alcance $2$) el escalón cae entre las distancias $2$ y $3$, y con núcleo $5$
(alcance $4$) entre $4$ y $5$. Sesenta celdas de sesenta, dos arquitecturas, tres semillas cada una,
dos componentes de la consulta y cinco distancias.

La consecuencia en el comportamiento es una falla específica y frecuente. En nuestra tarea las
preguntas tienen la forma \emph{cuál es la \textbf{<relación>} de \textbf{<entidad>}}, con la entidad
a distancia $1$ de la posición de lectura y la relación a distancia $3$. Con núcleo $3$ la relación
queda afuera de la ventana en el $100\,\%$ de las consultas, así que el modelo recupera solamente por
la entidad. Cuando se le pregunta por una relación que nunca se enunció, sobre una entidad que sí
apareció, responde con el hecho archivado equivocado en vez de abstenerse, y la abstención correcta
queda en $0{,}5850$ a $0{,}7349$ sobre tres semillas. Ensanchar el núcleo a $5$, que cuesta $1{,}280$
parámetros sobre $865{,}395$, la lleva a $\mathbf{0{,}9931}$ hasta $\mathbf{1{,}0000}$ sin solape
entre condiciones, y la exactitud global a $0{,}988$ hasta $0{,}993$ contra un piso trivial de
$0{,}4065$.

Después verificamos el mecanismo afuera de nuestro propio modelo. En \texttt{mamba-130m-hf}, de $129$
millones de parámetros, cambiar un token de cinco a ocho posiciones antes de la lectura mueve la
salida de la capa en todas las distancias y deja la salida de la convolución en \textbf{cero exacto},
$80$ celdas de $80$. El estado ve la secuencia entera y la consulta que lo lee no. De paso, el tap más
viejo de la convolución en ese checkpoint vale cero exacto en las $24$ capas, así que su alcance
efectivo es $2$ y no el nominal $3$. Informamos la medición y no una causa.

Damos el diagnóstico como una receta que no necesita entrenamiento. Se cambia un token de la consulta
y se mira si la recuperación se mueve. Si no se mueve, ese token es invisible para la búsqueda, y la
confianza que el modelo muestre sobre él no está ganada.
\end{abstract}

\section{El problema}

Un modelo con memoria persistente tiene que hacer dos cosas que suelen medirse juntas y en realidad
son separadas. Tiene que \emph{encontrar} la entrada correcta, y tiene que \emph{saber} cuándo la
entrada que busca no está. La segunda es la que separa una memoria de un fabulador, y es la difícil de
las dos.

Encontramos que en nuestro propio sistema la segunda falla tenía una causa puramente mecánica, aguas
arriba de todo lo que el modelo aprende. La consulta con la que se busca en el archivo se arma con una
convolución causal corta sobre los estados ocultos. En una pregunta de la forma \emph{cuál es
\texttt{<art>} \texttt{<sust>} de \texttt{<entidad>}}, la entidad queda un token antes de la posición
de lectura y el sustantivo que nombra la relación queda tres tokens antes. Un núcleo de tamaño $3$
alcanza dos tokens hacia atrás. La relación cae un token afuera de la ventana, de forma determinista,
en todas y cada una de las consultas.

El modelo no estaba ignorando la relación. No la podía ver.

\section{Trabajo relacionado}

\textbf{Las convoluciones cortas están en todos lados y su tamaño se hereda, no se mide.} Mamba-2
\citep{mamba2} usa por omisión una convolución causal depthwise de tamaño $4$. En la implementación de
referencia de Qwen3-Next, que usa Gated DeltaNet \citep{gateddeltanet} en sus capas de atención lineal,
la configuración expone \texttt{linear\_conv\_kernel\_dim}, documentado como el tamaño de núcleo de la
convolución de las capas de atención lineal y \textbf{con valor por omisión 4} \citep{qwen3next}. Las
ablaciones de esta línea informan que sacar la convolución empeora el modelo, lo cual establece que el
contexto local importa. Ninguna informa qué se pierde cuando la ventana está presente pero es
demasiado angosta para la consulta que hay que armar. Un alcance de tres tokens no es una
configuración exótica, es lo que corre en modelos desplegados.

\textbf{El resultado teórico más cercano.} \citet{cat} estudian el recuerdo asociativo de $n$-gramas y
prueban, por construcción, que alcanza con un filtro causal de largo $N$, igualando el filtro al
$n$-grama. Es la misma familia de idea y la citamos como precursora. Tres cosas la separan de lo que
informamos acá. Su consulta es un $n$-grama \emph{contiguo}, los últimos $N$ tokens sin nada
irrelevante en el medio, mientras que lo que manda en nuestro caso es la \emph{distancia} de un
componente, con relleno entre las dos partes que importan. Su Teorema 1 es un resultado de existencia
y no corren ningún experimento con el filtro demasiado corto. Y en la práctica informan lo contrario,
al decir que la convolución sobre consultas y claves ayuda en las construcciones teóricas pero no da
beneficios apreciables en los experimentos reales. Nosotros medimos un efecto grande en un régimen que
ellos no prueban, que es una memoria persistente \emph{entre secuencias} con una consulta compuesta.

\textbf{Por qué el campo no lo pudo ver.} La tarea canónica para esta capacidad, MQAR
\citep{zoology}, tiene una consulta de \emph{un solo token}. Se busca la clave y se devuelve el valor.
Por construcción esa consulta no puede tener una parte adentro de la ventana y otra afuera. Zoology
atribuye el $82\,\%$ de la brecha de recuerdo a que la convolución aplica un filtro fijo e
independiente de la entrada, que es una afirmación sobre \emph{adaptabilidad}. La nuestra es sobre
\emph{alcance}, y las dos son independientes. Nuestro modelo de núcleo $5$ tiene exactamente la misma
independencia de la entrada y la falla desaparece.

\textbf{Del lado de la atención.} La atención multi-token \citep{mta} convoluciona sobre los pesos de
atención para que la atención pueda condicionarse en varios tokens a la vez, motivada por la
observación de que un solo token no alcanza para localizar una coincidencia. Va en la misma dirección
y confirma la nuestra por contraste. Nunca se pregunta qué parte de la consulta queda ciega, y no
conecta la pregunta con la abstención.

\section{Montaje experimental}

El sustrato es un modelo de lenguaje de $3{,}5$\,MB, $865{,}395$ parámetros, $d=128$, cuatro bloques,
entrenado desde cero sobre un idioma sintético cerrado de $242$ tokens con un archivo persistente
co-entrenado. Los hechos se enuncian en una secuencia, se revisan en otra y se consultan en una
tercera, \textbf{con el estado recurrente reseteado entre secuencias}, así que la recuperación no se
puede resolver arrastrando activaciones hacia adelante.

Las consultas se leen en el último token de la pregunta. La lectura se inyecta en el bloque $0$
\textbf{antes del mixer}, y eso importa para el mecanismo. En ese punto el estado oculto todavía es la
embedding del token, sin ninguna mezcla recurrente, así que la ventana de la convolución es
literalmente todo lo que la consulta puede ver. Eso también explica por qué el efecto de abajo puede
ser cero exacto y no apenas chico.

Se puntúan por separado cuatro categorías de respuesta. \texttt{vigente} y \texttt{anterior} son
preguntas con respuesta en el archivo. \texttt{nose\_ent} pregunta por una entidad que nunca se
mencionó, que es la ausencia fácil. \texttt{nose\_rel} pregunta por una relación que nunca se enunció
\emph{sobre una entidad que sí apareció}, que es el caso que se parece a una alucinación real, porque
la recuperación encuentra un hecho vecino y genuino.

\section{Resultado 1. La ventana es una función escalón, y afuera vale cero exacto}

Medimos la sensibilidad reemplazando un componente de la pregunta por otro del mismo tipo y
registrando la distancia en variación total entre las distribuciones de lectura. Para variar la
distancia sin tocar la pregunta, la posición de lectura se corre $r$ lugares sobre el relleno que ya
sigue al signo de pregunta, que es exactamente equivalente a agregar $r$ rellenos.

\begin{center}
\begin{tabular}{lccccccc}
\toprule
modelo & núcleo (alcance) & $d{=}1$ & $d{=}2$ & $d{=}3$ & $d{=}4$ & $d{=}5$ & $d{=}6$ \\
\midrule
\texttt{v3} (3 semillas) & 3 (2) & $0{,}71$--$0{,}74$ & $0{,}15$--$0{,}25$ & $\mathbf{0{,}000000}$ & $\mathbf{0{,}000000}$ & $\mathbf{0{,}000000}$ & $\mathbf{0{,}000000}$ \\
\texttt{kq3} (3 semillas) & 5 (4) & $0{,}36$--$0{,}46$ & $0{,}10$--$0{,}15$ & $0{,}05$--$0{,}35$ & $0{,}07$--$0{,}33$ & $\mathbf{0{,}000000}$ & $\mathbf{0{,}000000}$ \\
\bottomrule
\end{tabular}
\end{center}

Sesenta celdas de sesenta coinciden con la predicción hecha antes de correr, que la sensibilidad es
positiva mientras $d \le$ alcance y cero más allá. Una celda daba al principio $0{,}0106$ donde iba un
cero. La causa era que la formulación por el valor anterior es dos tokens más larga, así que correr la
posición de lectura la recortaba contra el borde del tensor, y recortar \emph{acorta} la distancia que
se está midiendo. Esas muestras ahora se descartan en vez de recortarse y la celda pasó a
$0{,}000000$.

\section{Resultado 2. Ensanchar la ventana compra abstención}

\begin{center}
\begin{tabular}{lccccc}
\toprule
unidad & núcleo & \texttt{vigente} & \texttt{nose\_ent} & \textbf{\texttt{nose\_rel}} & abstención falsa \\
\midrule
\texttt{kq3\_s0} & 5 & $1{,}0000$ & $0{,}9538$ & $\mathbf{0{,}9931}$ & $0{,}0000$ \\
\texttt{kq3\_s1} & 5 & $0{,}9964$ & $0{,}9726$ & $\mathbf{1{,}0000}$ & $0{,}0030$ \\
\texttt{kq3\_s2} & 5 & $1{,}0000$ & $0{,}9356$ & $\mathbf{1{,}0000}$ & $0{,}0000$ \\
\midrule
\texttt{v3\_s0} & 3 & $1{,}0000$ & $0{,}9851$ & $0{,}6090$ & $0{,}0000$ \\
\texttt{v3\_s1} & 3 & $1{,}0000$ & $1{,}0000$ & $0{,}5850$ & $0{,}0000$ \\
\texttt{v3\_s2} & 3 & $0{,}9927$ & $0{,}9414$ & $0{,}7349$ & $0{,}0064$ \\
\bottomrule
\end{tabular}
\end{center}

No hay solape entre condiciones en \texttt{nose\_rel}. La exactitud global, contando una abstención
correcta como acierto, va de $0{,}988$ a $0{,}993$ con núcleo $5$ contra un piso trivial de
$0{,}4065$, que es lo que saca un modelo que se abstiene en todo.

\textbf{El costo existe y lo declaramos.} \texttt{nose\_ent}, el caso fácil, baja un poco, de
$0{,}941$--$1{,}000$ a $0{,}936$--$0{,}973$. Una lectura que informe solamente un número agregado de
abstención esconde ese intercambio.

\textbf{No es capacidad.} El modelo de núcleo $5$ agrega $1{,}280$ parámetros, el $0{,}15\,\%$, y el
control de núcleo $3$ ya recupera la entrada correcta con probabilidad $1{,}0000$. Lo que cambia es el
acceso, y el Resultado 1 lo mide directo.

\section{Resultado 3. Apagar el tap que alcanza la relación}

Apagar de a un tap por vez en el modelo de núcleo $5$ ya entrenado, sin reentrenar nada, daña
\texttt{vigente} en $0{,}483$, $0{,}582$ y $0{,}530$ para el tap que alcanza la relación, contra
$0{,}372$, $0{,}013$ y $0{,}195$ para el tap que alcanza una palabra funcional y $0{,}218$, $0{,}093$
y $0{,}157$ para el que alcanza el artículo. El tap de la relación tiene la magnitud aprendida
\emph{más chica} de los cinco, así que por unidad de peso cuesta cerca del doble que cualquiera de los
dos controles.

Un criterio preregistrado formulado sobre la métrica de abstención falló, e informamos por qué. Apagar
el tap saca la relación de \emph{todas} las consultas, así que el modelo se abstiene más, y una
métrica que premia abstenerse no puede bajar. El daño aparece en la exactitud y en la abstención
falsa.

\section{Resultado 4. La disociación se sostiene en un modelo real}

Todo lo anterior está medido en un modelo chico, así que verificamos el mecanismo en uno que no es
nuestro. \texttt{state-spaces/mamba-130m-hf}, $129{,}135{,}360$ parámetros, en CPU y sin entrenar
nada. En Mamba la convolución corta se aplica a $x$ \emph{antes} de calcular $B$, $C$ y $\Delta$, y
$C_t$ es el análogo de la consulta de lectura.

Cambiamos un solo token a distancia $d$ de una posición de lectura y medimos, en la capa $0$, el
movimiento de la salida de la convolución y el de la salida de la capa. Diez posiciones de lectura,
dos textos.

\begin{center}
\begin{tabular}{lcccccc}
\toprule
& $d{=}1$ & $d{=}2$ & $d{=}3$ & $d{=}5$ & $d{=}7$ & $d{=}8$ \\
\midrule
salida de \texttt{conv1d} & $0{,}8$--$3{,}7$ & $0{,}5$--$1{,}6$ & $\mathbf{0{,}0}$ & $\mathbf{0{,}0}$ & $\mathbf{0{,}0}$ & $\mathbf{0{,}0}$ \\
salida de la capa & $9\!\times\!10^{-2}$ & $2\!\times\!10^{-2}$ & $2\!\times\!10^{-2}$ & $1\!\times\!10^{-2}$ & $9\!\times\!10^{-2}$ & $1\!\times\!10^{-2}$ \\
\bottomrule
\end{tabular}
\end{center}

La salida de la capa se mueve en \emph{todas} las distancias, $80$ celdas de $80$, así que el estado
sí ve la secuencia entera. La salida de la convolución vale cero exacto más allá de su alcance.
\textbf{Eso es la disociación.} El modelo ve el contexto y la consulta que lee su estado no lo ve.

\textbf{Un hallazgo incidental, informado con su incertidumbre.} Habíamos predicho movimiento en
$d=3$, porque el núcleo es $4$. Vale cero exacto. La razón es que \textbf{el tap más viejo de la
convolución vale cero exacto en las $24$ capas}, lo que confirmamos por los pesos y también por
intervención, ya que cambiar un token mueve la salida de la convolución en exactamente tres posiciones
consecutivas y no cuatro. La ventana efectiva de este checkpoint es $3$ y no $4$, y el alcance es $2$.
No sabemos la causa. Un peso que vale cero exacto en $24 \times 1536$ entradas no sale de un
gradiente, así que un artefacto de conversión es plausible, pero no lo verificamos y no lo afirmamos.
El hecho se reproduce en tres líneas. Su consecuencia práctica, si generaliza, es que el núcleo
nominal es una cota superior del alcance y conviene medirlo en vez de suponerlo.

\section{El diagnóstico, como receta}

La medición del Resultado 1 no necesita entrenamiento ni gradientes. Para cualquier modelo que
consulte una memoria con una consulta armada desde una ventana local.

\begin{enumerate}
\item Tomar una consulta con sus partes separadas, de modo que algún componente quede a distancia $d$
de la posición de lectura.
\item Reemplazar ese componente por otro del mismo tipo.
\item Comparar las distribuciones de recuperación.
\end{enumerate}

Si son idénticas, ese componente es invisible para la búsqueda, y la confianza que el modelo muestre
sobre él no está ganada. Con un núcleo de $4$, que es el valor por omisión en las capas de atención
lineal desplegadas, el alcance es de tres tokens, así que cualquier parte de una pregunta que quede
más atrás no entra en la consulta de esa capa.

\section{Limitaciones}

El sustrato es un modelo chico sobre un idioma sintético cerrado, y la distancia entre la relación y
la posición de lectura la fija el generador, así que el efecto acá es determinista donde en texto
natural sería una distribución. La medición es de \emph{acceso} y no de uso. Un token adentro de la
ventana no necesariamente se usa, y de hecho el tap del artículo está adentro y aporta poco. Por
último, la lectura en nuestro modelo se inyecta antes del mixer, que es lo que hace que el cero sea
exacto. En un modelo profundo que consulta su memoria desde capas más altas, la recurrencia ya movió
información hacia adelante y el alcance efectivo crece con la profundidad. La afirmación se hace para
memoria consultada desde capas tempranas, que es donde vive.

\begin{thebibliography}{99}
\bibitem[Dao y Gu(2024)]{mamba2} Dao, T., y Gu, A. (2024). Transformers are SSMs. Generalized Models and Efficient Algorithms Through Structured State Space Duality. ICML 2024.
\bibitem[Yang et al.(2025)]{gateddeltanet} Yang, S., Kautz, J., y Hatamizadeh, A. (2025). Gated Delta Networks. Improving Mamba2 with Delta Rule. ICLR 2025. arXiv:2412.06464.
\bibitem[Hugging Face(2026)]{qwen3next} Documentación de Hugging Face Transformers, \texttt{Qwen3NextConfig}, parámetro \texttt{linear\_conv\_kernel\_dim}, valor por omisión $4$. Consultada el 2 de septiembre de 2026.
\bibitem[Li et al.(2024)]{cat} Li, M., Zhang, X., Huang, Y., y Oymak, S. (2024). On the Power of Convolution Augmented Transformer. arXiv:2407.05591.
\bibitem[Arora et al.(2024)]{zoology} Arora, S., Eyuboglu, S., Timalsina, A., Johnson, I., Poli, M., Zou, J., Rudra, A., y R\'e, C. (2024). Zoology. Measuring and Improving Recall in Efficient Language Models. ICLR 2024. arXiv:2312.04927.
\bibitem[Golovneva et al.(2025)]{mta} Golovneva, O., Wang, T., Weston, J., y Sukhbaatar, S. (2025). Multi-Token Attention. arXiv:2504.00927.
\end{thebibliography}

\end{document}
